% ch9-summary.tex
\chapter{Summary and Navigation}
\label{ch:summary}

This chapter answers the six questions the companion set out to answer,
each with pointers to the sections that show it, and ends with the
navigation table from ulm.lam examples to the Field Guide.

% ---------------------------------------------------------------------------
\section{Six answers}
\label{sec:summary-answers}

\qparagraph{1. Which guarantees do ulm.lam's types express?} Static sizes
(checked, \cref{sec:shapes-static}); the element type and rank (checked);
a conjunction of up to fifteen matrix facts in the tensor encoding, closed
under implication, with no contradictions and consistent with a static
shape (checked, \cref{sec:infra-props}); and, inside a function, that every
fact on a result follows from the facts on the operands by ulm.lam's rules
(checked, \cref{sec:infra-compat}). The facts themselves are guarantees
only relative to the declarations at the parameters and in \code{assume}.

\qparagraph{2. Which facts are assumed, verified, inferred or checked at run
time?} Assumed: the properties of parameters and of \code{assume}, and full
\emph{column} rank for dynamic sizes (\cref{sec:facts-declared},
\cref{sec:gls-infer}). Verified: shapes, contradictions, result types
(\cref{ch:infra}). Inferred: result properties (\cref{sec:facts-inferred}).
Established by analysis: sizes of dynamic dimensions, reuse under
dominance, in-place writes (\cref{sec:facts-analysis}). Checked at run
time: array sizes (Julia/Python runtime), positive definiteness
(\code{potrf}), exact singularity (\code{getrf}), workspace size
(\cref{sec:facts-runtime}). Not established: symmetry of inputs, \(n\ge m\),
zeros in the unreferenced triangle, conditioning.

\qparagraph{3. What does MLIR provide, and what does ulm.lam implement?}
MLIR: the tensor and memref types with an encoding slot, ODS, verifier
hooks, \code{InferTypeOpInterface}, the effect interfaces and generic passes
(CSE, DCE, LICM, canonicalization), pattern rewriting, dialect conversion,
One-Shot Bufferize, EmitC. ulm.lam: the property attribute and its meaning,
the verifiers, the inference rules, the effect declarations, the normal
form, kernel selection with its cost model, the DPS kernel ops and their
bufferization models, the workspace and \code{info} passes, the lowering
patterns, and the frontends (\cref{tab:infra-who}).

\qparagraph{4. How do those facts enable transformations?} \code{symmetric}
removes transposes (\cref{sec:gls-normal}). \code{spd} selects Cholesky
instead of LU, and with it \code{syrk} instead of \code{gemm} and two
triangular solves instead of an LU solve (\cref{tab:variants}).
\code{full_rank} turns \(X\T M\inv X\) from SPSD into SPD and so decides
the factorization of \(S\). Triangularity removes a factorization. Purity
lets LICM hoist an inverse out of a loop (\cref{sec:effects-licm}).
Read-only intents decide copies (\cref{sec:storage-intent}).

\qparagraph{5. How do value semantics, effects, aliasing and bufferization
interact?} All optimization happens on tensors, where values have no
identity and no aliasing, and effects are absent: uses can be shared,
pure ops moved and dead ones removed, including a dead factorization with
its check (\cref{ch:values}). DPS kernels name the value they overwrite;
One-Shot Bufferize turns each into an in-place write or a copy, honoring
read-only parameters (\cref{ch:storage}). After bufferization the kernels
have read/write effects on one default resource, no further optimization
needs alias information, and the workspace pass places temporaries at
disjoint offsets. Failure travels as a value (\code{info}) until it becomes
control flow (\cref{sec:effects-info}).

\qparagraph{6. Where is the corresponding theory?} In the table below.

% ---------------------------------------------------------------------------
\section{Navigation table}
\label{sec:summary-nav}

\subsection*{From Julia to typed IR}
{\small
\begin{tabularx}{\linewidth}{@{}P{3.6cm}P{4.0cm}LP{2.0cm}@{}}
\toprule
\textbf{ulm.lam example} & \textbf{MLIR mechanism} & \textbf{Theoretical concept} & \textbf{Field Guide} \\
\midrule
\code{@lam} traced into \code{ulm_lam} ops (\cref{sec:gls-trace}) & textual IR from overloaded operators; \code{func.func} & tracing front end of an embedded DSL & \FGs{3.7} \\
Julia \code{if} and loops disappear at trace time (\cref{sec:gls-trace}) & --- & binding time, staging & \FGs{8.6} \\
\bottomrule
\end{tabularx}
}

\subsection*{Facts: refinement, inference, assumptions, checks}
{\small
\begin{tabularx}{\linewidth}{@{}P{3.6cm}P{4.0cm}LP{2.0cm}@{}}
\toprule
\textbf{ulm.lam example} & \textbf{MLIR mechanism} & \textbf{Theoretical concept} & \textbf{Field Guide} \\
\midrule
\code{tensor<?x?xf64, #ulm_lam.props<spd>>} (\cref{sec:gls-types}) & attribute in the builtin tensor encoding; \code{VerifiableTensorEncoding} & refinement type & \FGs{1.7} \\
\code{spd} on a \(3\times4\) tensor rejected (\cref{sec:infra-props}) & \code{verifyEncoding} & well-formedness of refinements & \FGs{1.7} \\
result declared \code{spd} rejected without \code{full_rank} (\cref{sec:infra-compat}) & \code{InferTypeOpInterface}, \code{isCompatibleReturnTypes} & subtyping by implication; checking mode & \FGs{1.7}, \FGs{1.10} \\
\(X\T M\inv X\) inferred SPD (\cref{sec:gls-infer}) & \code{inferReturnTypes}, pass \code{ulm-lam-infer-properties} & type synthesis & \FGs{1.10} \\
property sets, closure, ``absence = unknown'' (\cref{sec:facts-lattice}) & bit-set attribute, \code{closure} & lattice, abstract domain, transfer functions & \FGs{5.1}, \FGs{5.2} \\
loop-carried facts re-declared with \code{assume} (\cref{sec:facts-lattice}) & block arguments keep their types & fixpoint over loops (not implemented) & \FGs{5.1} \\
\code{assume(S, :spd)} (\cref{sec:facts-declared}) & \code{ulm_lam.assume}, no inference hook & trusted cast into a refinement & \FGs{1.7} \\
\code{potrf} returns \code{info} (\cref{sec:effects-info}) & result value, then \code{scf.if} from \code{ulm-lam-info} & run-time check; status passing vs exceptions & \FGs{1.6} \\
\bottomrule
\end{tabularx}
}

\subsection*{Dimensions}
{\small
\begin{tabularx}{\linewidth}{@{}P{3.6cm}P{4.0cm}LP{2.0cm}@{}}
\toprule
\textbf{ulm.lam example} & \textbf{MLIR mechanism} & \textbf{Theoretical concept} & \textbf{Field Guide} \\
\midrule
inner sizes of a product (\cref{sec:shapes-static}) & op verifiers on static shapes & sizes in types & \FGs{1.8} \\
\code{Mat(:n, :m)} and \code{ulm_lam.dims} (\cref{sec:shapes-dims}) & \code{index} parameters, argument attribute, \code{dimOf} & named symbolic dimensions; indexed types (by analogy) & \FGs{10.4}, \FGs{1.8}, \FGs{6.1} \\
sizes bound at the call (\cref{sec:shapes-runtime}) & --- (Julia runtime) & run-time shape check & \FGs{10.4} \\
\bottomrule
\end{tabularx}
}

\subsection*{Values, uses and effects}
{\small
\begin{tabularx}{\linewidth}{@{}P{3.6cm}P{4.0cm}LP{2.0cm}@{}}
\toprule
\textbf{ulm.lam example} & \textbf{MLIR mechanism} & \textbf{Theoretical concept} & \textbf{Field Guide} \\
\midrule
\(W\) used twice, stored once (\cref{sec:values-uses}) & SSA use lists; \code{hasOneUse} in patterns & contraction and weakening; value semantics & \FGs{1.1}, \FGs{2.7} \\
\code{inverse} hoisted out of a loop (\cref{sec:effects-licm}) & \code{Pure}, LICM & effects decide code motion & \FGs{1.6}, \FGs{3.2} \\
factorization fails in a zero-trip loop (\cref{sec:effects-speculation}) & \code{AlwaysSpeculatable} (part of \code{Pure}) & speculation; honesty of effect declarations & \FGs{1.6}, \FGs{3.2} \\
unused tensor \code{potrf} erased (\cref{sec:effects}) & \code{MemoryEffectOpInterface}, DCE & effect-free ops are removable & \FGs{1.6} \\
\code{ulm_lamk.call} never moved (\cref{sec:effects}) & read/write of the default resource & worst-case effects of unknown calls & \FGs{1.6} \\
\code{scf.for}, \code{scf.if} bodies (\cref{sec:values-regions}) & MLIR regions & structured control flow (not region memory) & \FGs{2.6}, \FGs{1.5} \\
\bottomrule
\end{tabularx}
}

\subsection*{Storage}
{\small
\begin{tabularx}{\linewidth}{@{}P{3.6cm}P{4.0cm}LP{2.0cm}@{}}
\toprule
\textbf{ulm.lam example} & \textbf{MLIR mechanism} & \textbf{Theoretical concept} & \textbf{Field Guide} \\
\midrule
\code{trsm ... outs(\%B)} on tensors and memrefs (\cref{sec:storage-dps}) & \code{DestinationStyleOpInterface}, \code{DstBufferizableOpInterfaceExternalModel} & destination-passing style & \FGs{10.3} \\
copies of \(M\), \(X\), \(y\); none with \code{InOut} (\cref{sec:storage-intent}) & One-Shot Bufferize, \code{bufferization.writable} & in-place condition; uniqueness (by analogy) & \FGs{10.3}, \FGs{1.4} \\
loop state in one buffer (\cref{sec:storage-loops}) & \code{materialize_in_destination} & bufferization of loop-carried values & \FGs{10.3} \\
one workspace, 64-byte aligned views (\cref{sec:storage-workspace}) & pass \code{ulm-lam-workspace}, \code{memref.view} & arena allocation; buffer assignment & \FGs{10.3}, \FGs{4.1} \\
no overlap between kernel operands (\cref{sec:storage-alias}) & disjoint views; effects on \code{DefaultResource} & alias analysis (avoided) & \FGs{5.5}, \FGs{2.7} \\
properties dropped after selection (\cref{sec:storage-lost}) & \code{ulm-lam-strip-properties} & progressive lowering & \FGs{3.3} \\
\bottomrule
\end{tabularx}
}

\subsection*{Rewriting, kernels and floating point}
{\small
\begin{tabularx}{\linewidth}{@{}P{3.6cm}P{4.0cm}LP{2.0cm}@{}}
\toprule
\textbf{ulm.lam example} & \textbf{MLIR mechanism} & \textbf{Theoretical concept} & \textbf{Field Guide} \\
\midrule
normal form (\cref{sec:kernels-normal}) & \code{OpRewritePattern}, greedy driver & rewriting to a fixpoint; termination & \FGs{3.4}, \FGs{3.5} \\
\(X\T M\inv X\) as \code{trsm} + \code{syrk} (\cref{sec:gls-select}) & custom pass, dynamic program & instruction selection by tiling; matrix chain & \FGs{4.5} \\
\(W\) reused transposed (\cref{sec:kernels-select}) & memo table, dominance & CSE up to transposition; DAG covering & \FGs{4.5} \\
C++ from memrefs (\cref{sec:kernels-emitc}) & \code{ConversionTarget}, \code{TypeConverter}, \code{applyPartialConversion} & dialect conversion; legalization, ABI & \FGs{3.3}, \FGs{4.6} \\
last-bit differences; NaN-safe accumulation (\cref{sec:kernels-fp}) & --- & real vs floating-point semantics & --- \\
Cholesky fails for nearly rank-deficient \(X\) (\cref{sec:kernels-fp}) & --- & conditioning of normal equations & --- \\
\bottomrule
\end{tabularx}
}

\subsection*{Evaluation order and specialization}
{\small
\begin{tabularx}{\linewidth}{@{}P{3.6cm}P{4.0cm}LP{2.0cm}@{}}
\toprule
\textbf{ulm.lam example} & \textbf{MLIR mechanism} & \textbf{Theoretical concept} & \textbf{Field Guide} \\
\midrule
flop counts as costs (\cref{sec:orders-cost}) & inline formulas; \code{KernelCostOpInterface} unused & analytical cost model & \FGs{10.5} \\
\((QK\T)V\) or \(Q(K\T V)\) at run time (\cref{sec:orders-runtime}) & \code{scf.if} on \code{index} cost polynomials & symbolic parameters; residual program & \FGs{6.1}, \FGs{8.5} \\
operator hoisted if the trip count pays (\cref{sec:orders-hoist}) & \code{scf.if} before and in the loop & run-time decision on a symbolic parameter & \FGs{6.1} \\
\code{Float32} library (\cref{sec:orders-spec}) & same IR text with \code{f32} & monomorphization & \FGs{8.7} \\
\bottomrule
\end{tabularx}
}


The entries marked --- in the last column have no entry of their own in
the Field Guide; floating-point behavior is covered by the numerical
literature~\cite{higham}.
